\section[Cote 104 : catégories de modèles, trois exemples, un complexe globulaire]{Cote 104 : catégories de modèles, trois exemples, un complexe globulaire\protect\verifie{Folder104}}\label{sec:104}

La cote 104, \emph{[Champs (stacks) 1] : copies de notes manuscrites (s.d.)}, seize pages sans date, appartient au groupe \q{Autour de \emph{Pursuing stacks}}, que l'inventaire date \q{1968-1984}. Les pages~6 à~16, sous sa pagination~1 à~6, sont une rédaction suivie : une catégorie $M$ de \q{modèles} est remplacée par des ensembles ordonnés --- l'ensemble $\mathbb N$ de ses classes d'isomorphisme, les ensembles $I_n$ de sous-objets, les applications type $\tau_n$ --- dont elle se reconstruit. La page~6 pose les deux hypothèses --- dans la seconde, \q{automorphisme} est biffé et remplacé par \q{end.} --- et affirme que l'application $(*)$ ci-dessous \q{est injective par (i)(ii)} (\lot{104}{1}{6}) ; la page~7 ajoute \q{On a donc (par (ii)) $\tau_n^{-1}(n) = \{d_n\}$} (\lot{104}{1}{7}) ; la page~10 affirme la réciproque et donne trois exemples, avec leurs dénombrements en marge (\lot{104}{1}{10}). Parmi les feuillets d'essai qui précèdent, celui de la page~4 porte une suite $\mathbb Z[F_{n+1}] \to \mathbb Z[F_n] \xrightarrow{t_n - s_n} \mathbb Z[F_{n-1}]$ et, sous une accolade, le calcul $t_{n-1}(t_n y - s_n y) - s_{n-1}(t_n y - s_n y) = 0$ (\lot{104}{1}{4}).

\begin{definition}
Soit $M$ une catégorie telle que (i)~toute flèche de $M$ est un monomorphisme, et (ii)~tout endomorphisme d'un objet de $M$ est l'identité. Pour un objet $D$, un sous-objet de $D$ est une classe de monomorphismes de but $D$, deux monomorphismes $u, v$ étant équivalents s'ils se factorisent l'un par l'autre ; les sous-objets sont ordonnés par factorisation, et l'on note $u(D')$ le sous-objet défini par $u : D' \to D$. Pour un sous-objet $x$, on note encore $x$ un objet qui le représente. Les classes d'isomorphisme d'objets de $M$ sont préordonnées par l'existence d'une flèche.
\end{definition}

\begin{enonce}[Pages 6 et 7]
\begin{enumerate}
\item L'application $(*)$ : $\operatorname{Hom}_M(D', D) \to \widetilde D$, $u \mapsto u(D')$, est injective.
\item Si $f : D' \to D$ et $g : D \to D'$, alors $f$ est un isomorphisme d'inverse $g$ ; le préordre sur les classes est donc un ordre.
\item Si $x \leq y$ sont deux sous-objets de $D$ et s'il existe une flèche $y \to x$ --- en particulier si $x$ et $y$ sont de même type ---, alors $x = y$. Par suite $\tau_n^{-1}(n) = \{d_n\}$, et $\tau_n$ est strictement croissante.
\end{enumerate}
\end{enonce}

\begin{enonce}[Page 10, la réciproque : un pas]
Soit $f$ une application bijective d'un ensemble dans lui-même telle que $f \circ f = f$. Alors $f = \mathrm{id}$. En particulier la loi de composition $u_{d_n} \circ u_{d_n} = u_{d_n}$ force $u_{d_n} = \mathrm{id}$.
\end{enonce}

\begin{enonce}[Page 10, les trois exemples]
\begin{enumerate}
\item La catégorie des $\Delta_n = \{0, \dots, n\}$ et des applications strictement croissantes vérifie (i) et (ii), et l'ensemble $\mathfrak P^*(\Delta_n)$ des parties non vides de $\Delta_n$ a $2^{n+1} - 1$ éléments.
\item Les faces de $[0,1]^n$, qui fixent certaines coordonnées à $0$ ou à $1$ et laissent les autres libres, sont au nombre de $\sum_{k} \binom nk 2^k = 3^n$.
\item Sur $I_n = (\Delta_{n-1} \times \{\pm 1\}) \sqcup \{d_n\}$, la relation \q{$(k, \varepsilon) \leq (k', \varepsilon')$ si et seulement si $(k, \varepsilon) = (k', \varepsilon')$ ou $k < k'$, et $d_n$ plus grand élément} est un ordre ; $I_n$ a $2n + 1$ éléments, dont exactement deux de type $k$ pour chaque $k < n$. Avec $\{\pm 1\}$ muni de l'ordre chaotique, le produit n'est qu'un préordre : $(k, +1)$ et $(k, -1)$, distincts, sont chacun inférieurs à l'autre.
\end{enumerate}
\end{enonce}

\begin{enonce}[Page 4, le complexe globulaire]
Soient des applications $s_n, t_n : F_n \to F_{n-1}$ et $s_{n-1}, t_{n-1} : F_{n-1} \to F_{n-2}$ telles que $s_{n-1}s_n = s_{n-1}t_n$ et $t_{n-1}s_n = t_{n-1}t_n$. Alors les morphismes $t_n - s_n : \mathbb Z[F_n] \to \mathbb Z[F_{n-1}]$ et $t_{n-1} - s_{n-1} : \mathbb Z[F_{n-1}] \to \mathbb Z[F_{n-2}]$ induits sur les groupes abéliens libres composent en zéro.
\end{enonce}

\begin{proof}[Démonstration de l'énoncé des pages 6 et 7]
(1) Soient $u, v : D' \to D$ tels que $u(D') = v(D')$. Comme $u(D') \leq v(D')$, il existe $\theta : D' \to D'$ avec $v\theta = u$. Par (ii), $\theta = \mathrm{id}$, donc $u = v$. L'hypothèse (i) ne sert qu'à faire de $u$ et $v$ des monomorphismes, c'est-à-dire à donner un sens à $u(D')$.

(2) Le composé $gf$ est un endomorphisme de $D'$, et $fg$ un endomorphisme de $D$ ; par (ii), ce sont des identités.

(3) Écrivons $x \leq y$ sous la forme $a_x = a_y \circ j$, où $a_x, a_y$ sont les monomorphismes de but $D$ qui représentent $x$ et $y$ et $j : x \to y$. Soit $g : y \to x$. Le composé $j \circ g$ est un endomorphisme de $y$, donc l'identité par (ii), et $a_y = a_y \circ j \circ g = a_x \circ g$ : $a_y$ se factorise par $a_x$, c'est-à-dire $y \leq x$, et $x = y$. Si $\tau_n(x) = n$, $x$ est isomorphe à $D_n$, qui représente le plus grand sous-objet $d_n$ ; on a $x \leq d_n$ et une flèche $d_n \to x$, donc $x = d_n$. Si $x < y$, il y a une flèche $x \to y$, donc $\tau_n(x) \leq \tau_n(y)$ ; l'égalité donnerait une flèche $y \to x$, donc $x = y$ ; ainsi $\tau_n(x) < \tau_n(y)$ dans l'ensemble ordonné des classes.
\end{proof}

\begin{proof}[Démonstration du pas de la réciproque]
Pour tout $x$, $f(f(x)) = f(x)$, et l'injectivité de $f$ donne $f(x) = x$.
\end{proof}

\begin{proof}[Démonstration des trois exemples]
(1) Une application strictement croissante est injective, donc simplifiable à gauche : si $fg = fh$, alors $g = h$ point par point. Une application strictement croissante de $\Delta_n$ dans lui-même est l'identité : c'est l'unique application strictement croissante de $\{0, \dots, n\}$ sur une partie à $n + 1$ éléments de $\Delta_n$, à savoir $\Delta_n$ tout entier, et l'identité en est une. Les parties de $\Delta_n$ sont au nombre de $2^{n+1}$, dont une vide.

(2) Une face est la donnée, pour chaque coordonnée, de la valeur $0$, de la valeur $1$, ou de rien : c'est une application de $\{1, \dots, n\}$ dans un ensemble à trois éléments, et il y en a $3^n$ --- ce que la marge appelle une section partielle de $\Delta_{n-1} \times \{\pm 1\} \to \Delta_{n-1}$. En comptant les faces selon le nombre $k$ des coordonnées fixées, on obtient $\binom nk 2^k$ faces pour chaque $k$, et $\sum_k \binom nk 2^k = (2 + 1)^n$ par la formule du binôme ; c'est la formule de la page~5.

(3) La relation est réflexive. Elle est transitive : $d_n$ est au-dessus de tout et au-dessous de lui seul, et sur $\Delta_{n-1} \times \{\pm 1\}$ la composée de deux relations \q{égal ou de premier indice strictement plus petit} en est une. Elle est antisymétrique : si $(k, \varepsilon) \leq (k', \varepsilon')$ et $(k', \varepsilon') \leq (k, \varepsilon)$ avec des éléments distincts, on aurait $k < k'$ et $k' < k$. Le cardinal est $2n + 1$, et les éléments de type $k$ sont $(k, +1)$ et $(k, -1)$. Pour l'ordre chaotique, $(k, \varepsilon) \leq (k', \varepsilon')$ équivaut à $k \leq k'$, et $(k, +1)$, $(k, -1)$ sont chacun inférieurs à l'autre.
\end{proof}

\begin{proof}[Démonstration de l'énoncé de la page 4]
Le passage de $F$ à $\mathbb Z[F]$ est fonctoriel : l'application induite par un composé est le composé des applications induites. Le composé $(t_{n-1} - s_{n-1})(t_n - s_n)$ vaut donc $\mathbb Z[t_{n-1}t_n] - \mathbb Z[t_{n-1}s_n] - \mathbb Z[s_{n-1}t_n] + \mathbb Z[s_{n-1}s_n]$, et les identités globulaires annulent les termes deux à deux.
\end{proof}

\begin{remarque}
L'exemple~3 est démontré sous la forme que la lecture substitue à la page : $\Delta_{n-1}$ au lieu de $\Delta_n$, qui donnerait $2n + 3$ éléments et non les $2n + 1$ de la marge ; l'ordre des cellules du disque au lieu du produit par un préordre chaotique, qui n'est pas antisymétrique.
\end{remarque}

\begin{remarque}
L'hypothèse~(i) n'intervient ni dans l'ordre sur les classes ni dans $u_{d_n} = \mathrm{id}$, que (ii) donne seule ; dans $(*)$, que la page dit injective \q{par (i)(ii)}, elle sert seulement à définir $u(D')$.
\end{remarque}

\begin{remarque}
L'exemple~1 montre que la catégorie des $\Delta_n$ et des injections croissantes n'a pas de dégénérescence. Pour les exemples, seul le calcul sur les ensembles $I_n$ est démontré ; l'identification des sous-objets des catégories de cubes et de disques à ces ensembles est celle de la lecture (non démontré ici).
\end{remarque}

\begin{enonce}[Page 6, note de la lecture : la bonne fondation]
Soit $M$ une catégorie vérifiant (ii). Pour deux objets $X, Y$, écrivons $X \prec Y$ s'il existe une flèche $X \to Y$ et aucune flèche $Y \to X$.
\begin{enumerate}
\item Toute flèche $X \to Y$ qui n'est pas un isomorphisme donne $X \prec Y$ ; si $M$ est squelettique, toute flèche qui n'est pas une identité le donne.
\item Si $M$ est squelettique et si $\prec$ est bien fondée, le rang ordinal $\deg$ pour $\prec$ est un degré que toute flèche autre qu'une identité augmente strictement, et $M$ est une catégorie directe : une structure de Reedy dont la classe descendante est formée des identités et la classe ascendante de toutes les flèches.
\item Réciproquement, s'il existe une application $\deg$ des objets dans un ensemble ordonné bien fondé que toute flèche autre qu'une identité augmente strictement, $\prec$ est bien fondée.
\item Les hypothèses (i) et (ii) ne suffisent pas : la catégorie définie par l'ensemble ordonné $\mathbb Z$ les vérifie, et $\prec$ n'y est pas bien fondée.
\end{enumerate}
\end{enonce}

\begin{proof}
(1) Si $g : Y \to X$, alors $gf$ et $fg$ sont des endomorphismes, donc des identités par (ii), et $f$ serait un isomorphisme. Si $M$ est squelettique et $f$ un isomorphisme, $X = Y$, et $f$, endomorphisme de $X$, est l'identité par (ii) ; une flèche qui n'est pas une identité n'est donc pas un isomorphisme.

(2) Le rang d'un objet $X$ est le plus petit ordinal strictement supérieur aux rangs des objets $Z \prec X$ ; il est défini par récurrence bien fondée, et $Z \prec X$ entraîne $\deg Z < \deg X$. Par (1), toute flèche $X \to Y$ autre qu'une identité donne $\deg X < \deg Y$. La condition sur la classe descendante est vide, puisqu'elle n'a que des identités. Enfin toute flèche $f$ s'écrit d'une seule façon comme une identité suivie d'une flèche : si $f = p \circ i$ avec $i$ une identité, l'objet intermédiaire est la source de $f$ et $p = f$.

(3) Si $X \prec Y$, une flèche $f : X \to Y$ n'est pas une identité, puisque $X = Y$ donnerait la flèche $\mathrm{id}_X$ de $Y$ vers $X$ ; donc $\deg X < \deg Y$, et $\prec$ est contenue dans l'image réciproque par $\deg$ d'un ordre bien fondé, qui est bien fondée.

(4) Entre deux objets de $\mathbb Z$ il y a au plus une flèche : (i) et (ii) sont immédiats. Pour tout $m$, $m - 1 \prec m$, et $\mathbb Z$ n'a donc pas d'élément minimal pour $\prec$.
\end{proof}

\begin{remarque}
La lecture énonce qu'une catégorie vérifiant (i) et (ii) est directe au sens de Reedy \q{dès que} l'ordre de $\mathbb N$ est bien fondé : l'exemple de $\mathbb Z$ montre que cette condition est bien une hypothèse supplémentaire, et non une conséquence de (i) et (ii). La structure de Reedy est celle que définit la bibliothèque mathlib. L'hypothèse (i) n'y sert pas ; elle dit seulement que la classe ascendante est formée de monomorphismes. La catégorie squelettique est celle que la page~9 appelle \q{rigide} ; sans elle, un isomorphisme entre deux objets distincts n'augmenterait aucun degré, et c'est l'ordre sur les classes, non sur les objets, qui est bien fondé.
\end{remarque}

\begin{enonce}[Page 10, exemple 1 : les identités simpliciales]
Notons $\delta_i : \Delta_n \to \Delta_{n+1}$, pour $0 \leq i \leq n + 1$, l'application strictement croissante qui évite $i$.
\begin{enumerate}
\item Pour $i \leq j$, $\delta_{j+1} \circ \delta_i = \delta_i \circ \delta_j$.
\item Toute application strictement croissante $f : \Delta_m \to \Delta_{n+1}$ avec $m \leq n$ s'écrit $\delta_i \circ g$, avec $g : \Delta_m \to \Delta_n$ strictement croissante.
\item La catégorie des $\Delta_n$ et des applications strictement croissantes est squelettique, toute flèche autre qu'une identité y augmente strictement $n$, et c'est une catégorie directe au sens de l'énoncé précédent.
\end{enumerate}
\end{enonce}

\begin{proof}
(1) Les deux membres envoient $k \in \Delta_n$ sur $k$ si $k < i$, sur $k + 1$ si $i \leq k < j$, et sur $k + 2$ si $k \geq j$ : c'est une vérification cas par cas sur les inégalités entre $i$, $j$ et $k$. C'est la première identité simpliciale, que mathlib démontre dans la catégorie simpliciale ; elle est refaite ici dans la catégorie sans dégénérescences.

(2) Comme $f$ est injective, une surjection de $\Delta_m$, qui a $m + 1 \leq n + 1$ éléments, sur $\Delta_{n+1}$, qui en a $n + 2$, est impossible : $f$ évite un $i$. Chaque $f(k)$, distinct de $i$, s'écrit $\delta_i(g(k))$ pour un unique $g(k)$, et $g$ est strictement croissante parce que $\delta_i$ l'est et reflète donc l'ordre.

(3) Une flèche $\Delta_m \to \Delta_n$ est injective, donc $m \leq n$. Si $\Delta_m \simeq \Delta_n$, on a $m \leq n$ et $n \leq m$. Si $m = n$, la flèche est un endomorphisme, donc l'identité (exemple~1 ci-dessus). Le degré $n$, à valeurs dans les entiers, vérifie donc l'hypothèse de la partie (3) de l'énoncé précédent, et la partie (2) s'applique.
\end{proof}

\begin{remarque}
Ni les identités simpliciales ni la factorisation par les faces ne sont sur la page ; elles justifient l'identification, faite par la lecture, des préfaisceaux sur cette catégorie aux ensembles semi-simpliciaux : par (2) et une récurrence sur $n - m$, toute flèche est un composé de faces, et (1) donne, en passant aux préfaisceaux, les relations $d_i d_j = d_{j-1} d_i$ pour $i < j$. La récurrence elle-même n'est pas formalisée ici, seul son pas l'est. Dans cet exemple $\mathbb N = \mathbf N$, comme la page~10 le dit des \q{cas que j'ai en tête}, et la bonne fondation est celle des entiers.
\end{remarque}

\noindent Ne sont pas démontrés ici : la décomposition de toute flèche de l'exemple~1 en composé de faces au-delà du pas ci-dessus ; l'équivalence de $M$ avec la catégorie $M_0$ reconstruite à partir de $(\mathbb N, I_n, \tau_n, \varphi_{n,n'})$ (page~9), et la réciproque de la page~10 au-delà du pas ci-dessus ; la sphéricité des réalisations $|I_n^*|$ et la régularité des $I_n^*$ dans les trois exemples ; la question des suites régulières de la page~11 ; les préfaisceaux admissibles et les ensembles ordonnés $M$-épinglés des pages~12 à~16\piste{104-pinned-posets-nonsingular-presheaves}, avec l'équivalence de la page~13 ; les conditions $(C_i)$ de la page~4 sur les équivalences de systèmes globulaires.
